% GENERATED FILE. Edit canonical lessons, then run npm run generate:books.
% canonical-source-hash: fnv1a64:b7b22ff56bbadc9f
% canonical-lessons: TA-C345-kataicolli, TA-C345-nitikkatai, TA-C345-villuppattu, TA-C345-kataikkalam, TA-C345-arankeru

\chapter{Telling a Tale}
\label{ch:ta-telling-a-tale}
\hlchaptermodality{345}

\begin{chapteropening}
\textbf{By the end of this chapter:} \emph{I can talk about storytellers and how stories are told.}

Complete the last lesson of chapter 345: I can talk about storytellers and how stories are told.
\end{chapteropening}

\section[\texorpdfstring{\ta{கதைசொல்லி} (kataicolli)}{கதைசொல்லி (kataicolli)}]{\ta{கதைசொல்லி} (kataicolli) — a storyteller}

\label{lesson:TA-C345-kataicolli}



\begin{quote}
Before the new one: say the Tamil for to curse, then the Tamil for to change shape.
\end{quote}

\subsection*{You'll want to know: \ta{கதைசொல்லி}}
\textbf{\ta{கதைசொல்லி}} — \emph{kataicolli} — \textquotedblleft{}a storyteller\textquotedblright{}.

\begin{quote}
\ta{அந்தக்} \ta{கதைசொல்லி} \ta{குழந்தைகளைச்} \ta{சிரிக்க} \ta{வைத்தார்}.
\end{quote}

\begin{quote}
\emph{antak kataicolli kuḻantaikaḷaic cirikka vaittār.}
\end{quote}

\begin{quote}
That storyteller made the children laugh.
\end{quote}

\ta{கதை} + \ta{சொல்லி}, one who tells.

\begin{cousinweb}
\textbf{\ta{கதைசொல்லி}} (\emph{kataicolli}) joins \ta{கதை}, \textquotedblleft{}story\textquotedblright{}, from Sanskrit \emph{kathā} (CDIAL 2704), with \ta{சொல்லி}, \textquotedblleft{}one who tells\textquotedblright{}, from native \ta{சொல்} (DEDR 2855).
\end{cousinweb}

\begin{grammarlens}[title={What you've built}]
One more word for telling what happened.
\end{grammarlens}

\subsection*{Your turn}
\begin{itemize}
\raggedright
  \item \emph{Say it:} \emph{kataicolli}
  \item \emph{Say it:} \emph{kataicolli}, once more
  \item \emph{Say it:} \emph{antak kataicolli kuḻantaikaḷaic cirikka vaittār.} — That storyteller made the children laugh.
  \item \emph{Recall:} say the Tamil for to curse, then the Tamil for to change shape, then say \emph{kataicolli} again
\end{itemize}

\begin{tcolorbox}[breakable,colback=teal!4,colframe=teal!35!black,title={Before you move on}]
What does \ta{கதைசொல்லி} mean? (\textquotedblleft{}A storyteller\textquotedblright{}.) Say it once more.
\end{tcolorbox}
\section[\texorpdfstring{\ta{நீதிக்கதை} (nītikkatai)}{நீதிக்கதை (nītikkatai)}]{\ta{நீதிக்கதை} (nītikkatai) — a moral tale; a fable}

\label{lesson:TA-C345-nitikkatai}



\begin{quote}
Before the new one: say the Tamil for to change shape, then the Tamil for a storyteller.
\end{quote}

\subsection*{You'll want to know: \ta{நீதிக்கதை}}
\textbf{\ta{நீதிக்கதை}} — \emph{nītikkatai} — \textquotedblleft{}a moral tale; a fable\textquotedblright{}.

\begin{quote}
\ta{ஒவ்வொரு} \ta{நீதிக்கதையும்} \ta{ஒரு} \ta{பாடம்} \ta{சொல்கிறது}.
\end{quote}

\begin{quote}
\emph{ovvoru nītikkataiyum oru pāṭam colkiṟatu.}
\end{quote}

\begin{quote}
Every moral tale teaches a lesson.
\end{quote}

A story for children that ends with a lesson.

\begin{cousinweb}
\textbf{\ta{நீதிக்கதை}} (\emph{nītikkatai}) is built from \ta{நீதி}, \textquotedblleft{}right conduct\textquotedblright{}, borrowed from Sanskrit \emph{nīti} (CDIAL 7546), with \ta{கதை}, \textquotedblleft{}story\textquotedblright{}, from Sanskrit \emph{kathā} (CDIAL 2704).
\end{cousinweb}

\begin{grammarlens}[title={What you've built}]
One more word for telling what happened.
\end{grammarlens}

\subsection*{Your turn}
\begin{itemize}
\raggedright
  \item \emph{Say it:} \emph{nītikkatai}
  \item \emph{Say it:} \emph{nītikkatai}, once more
  \item \emph{Say it:} \emph{ovvoru nītikkataiyum oru pāṭam colkiṟatu.} — Every moral tale teaches a lesson.
  \item \emph{Recall:} say the Tamil for to change shape, then the Tamil for a storyteller, then say \emph{nītikkatai} again
\end{itemize}

\begin{tcolorbox}[breakable,colback=teal!4,colframe=teal!35!black,title={Before you move on}]
What does \ta{நீதிக்கதை} mean? (\textquotedblleft{}A moral tale; a fable\textquotedblright{}.) Say it once more.
\end{tcolorbox}
\section[\texorpdfstring{\ta{வில்லுப்பாட்டு} (villuppāṭṭu)}{வில்லுப்பாட்டு (villuppāṭṭu)}]{\ta{வில்லுப்பாட்டு} (villuppāṭṭu) — a bow-song}

\label{lesson:TA-C345-villuppattu}



\begin{quote}
Before the new one: say the Tamil for a storyteller, then the Tamil for a moral tale.
\end{quote}

\subsection*{You'll want to know: \ta{வில்லுப்பாட்டு}}
\textbf{\ta{வில்லுப்பாட்டு}} — \emph{villuppāṭṭu} — \textquotedblleft{}a bow-song\textquotedblright{}.

\begin{quote}
\ta{திருவிழாவில்} \ta{வில்லுப்பாட்டு} \ta{நடந்தது}.
\end{quote}

\begin{quote}
\emph{tiruviḻāvil villuppāṭṭu naṭantatu.}
\end{quote}

\begin{quote}
There was a bow-song performance at the festival.
\end{quote}

A sung story, told with a large bow as the main instrument.

\begin{cousinweb}
\textbf{\ta{வில்லுப்பாட்டு}} (\emph{villuppāṭṭu}) is a Tamil compound built from \ta{வில்}, \textquotedblleft{}bow\textquotedblright{} (DEDR 5422), and \ta{பாட்டு}, \textquotedblleft{}song\textquotedblright{} (DEDR 4065).
\end{cousinweb}

\begin{grammarlens}[title={What you've built}]
One more word for telling what happened.
\end{grammarlens}

\subsection*{Your turn}
\begin{itemize}
\raggedright
  \item \emph{Say it:} \emph{villuppāṭṭu}
  \item \emph{Say it:} \emph{villuppāṭṭu}, once more
  \item \emph{Say it:} \emph{tiruviḻāvil villuppāṭṭu naṭantatu.} — There was a bow-song performance at the festival.
  \item \emph{Recall:} say the Tamil for a storyteller, then the Tamil for a moral tale, then say \emph{villuppāṭṭu} again
\end{itemize}

\begin{tcolorbox}[breakable,colback=teal!4,colframe=teal!35!black,title={Before you move on}]
What does \ta{வில்லுப்பாட்டு} mean? (\textquotedblleft{}A bow-song\textquotedblright{}.) Say it once more.
\end{tcolorbox}
\section[\texorpdfstring{\ta{கதைக்களம்} (kataikkaḷam)}{கதைக்களம் (kataikkaḷam)}]{\ta{கதைக்களம்} (kataikkaḷam) — a setting; a plot}

\label{lesson:TA-C345-kataikkalam}



\begin{quote}
Before the new one: say the Tamil for a moral tale, then the Tamil for a bow-song.
\end{quote}

\subsection*{You'll want to know: \ta{கதைக்களம்}}
\textbf{\ta{கதைக்களம்}} — \emph{kataikkaḷam} — \textquotedblleft{}a setting; a plot\textquotedblright{}.

\begin{quote}
\ta{இந்தப்} \ta{படத்தின்} \ta{கதைக்களம்} \ta{ஒரு} \ta{சிறிய} \ta{கிராமம்}.
\end{quote}

\begin{quote}
\emph{intap paṭattiṉ kataikkaḷam oru ciṟiya kirāmam.}
\end{quote}

\begin{quote}
The setting of this film is a small village.
\end{quote}

\ta{களம்}, field: the ground the story stands on.

\begin{cousinweb}
\textbf{\ta{கதைக்களம்}} (\emph{kataikkaḷam}) joins \ta{கதை}, \textquotedblleft{}story\textquotedblright{}, from Sanskrit \emph{kathā} (CDIAL 2704), with native \ta{களம்}, \textquotedblleft{}place, open space\textquotedblright{} (DEDR 1376).
\end{cousinweb}

\begin{grammarlens}[title={What you've built}]
One more word for telling what happened.
\end{grammarlens}

\subsection*{Your turn}
\begin{itemize}
\raggedright
  \item \emph{Say it:} \emph{kataikkaḷam}
  \item \emph{Say it:} \emph{kataikkaḷam}, once more
  \item \emph{Say it:} \emph{intap paṭattiṉ kataikkaḷam oru ciṟiya kirāmam.} — The setting of this film is a small village.
  \item \emph{Recall:} say the Tamil for a moral tale, then the Tamil for a bow-song, then say \emph{kataikkaḷam} again
\end{itemize}

\begin{tcolorbox}[breakable,colback=teal!4,colframe=teal!35!black,title={Before you move on}]
What does \ta{கதைக்களம்} mean? (\textquotedblleft{}A setting; a plot\textquotedblright{}.) Say it once more.
\end{tcolorbox}
\section[\texorpdfstring{\ta{அரங்கேறு} (araṅkēṟu)}{அரங்கேறு (araṅkēṟu)}]{\ta{அரங்கேறு} (araṅkēṟu) — to be staged}

\label{lesson:TA-C345-arankeru}



\begin{quote}
Before the new one: say the Tamil for a bow-song, then the Tamil for a setting.
\end{quote}

\subsection*{You'll want to know: \ta{அரங்கேறு}}
\textbf{\ta{அரங்கேறு}} — \emph{araṅkēṟu} — \textquotedblleft{}to be staged\textquotedblright{}.

\begin{quote}
\ta{அவளுடைய} \ta{நடனம்} \ta{நேற்று} \ta{அரங்கேறியது}.
\end{quote}

\begin{quote}
\emph{avaḷuṭaiya naṭaṉam nēṟṟu araṅkēṟiyatu.}
\end{quote}

\begin{quote}
Her dance was staged for the first time yesterday.
\end{quote}

Past (it …): \ta{அரங்கேறியது}, \emph{araṅkēṟiyatu}.

Also of a dancer's first public show.

\begin{cousinweb}
\textbf{\ta{அரங்கேறு}} (\emph{araṅkēṟu}) joins \ta{அரங்கு}, \textquotedblleft{}a stage, a hall\textquotedblright{}, which DEDR lists among its loans (DEDR App. 8), with native \ta{ஏறு}, \textquotedblleft{}to climb, go up\textquotedblright{} (DEDR 916).
\end{cousinweb}

\begin{grammarlens}[title={What you've built}]
One more word for telling what happened.
\end{grammarlens}

\subsection*{Your turn}
\begin{itemize}
\raggedright
  \item \emph{Say it:} \emph{araṅkēṟu}
  \item \emph{Say it:} \emph{araṅkēṟu}, once more
  \item \emph{Say it:} \emph{avaḷuṭaiya naṭaṉam nēṟṟu araṅkēṟiyatu.} — Her dance was staged for the first time yesterday.
  \item \emph{Recall:} say the Tamil for a bow-song, then the Tamil for a setting, then say \emph{araṅkēṟu} again
\end{itemize}

\begin{tcolorbox}[breakable,colback=teal!4,colframe=teal!35!black,title={Before you move on}]
What does \ta{அரங்கேறு} mean? (\textquotedblleft{}To be staged\textquotedblright{}.) Say it once more.
\end{tcolorbox}
